
\documentclass{article}
\usepackage{colortbl}
\usepackage{makecell}
\usepackage{multirow}
\usepackage{supertabular}

\begin{document}

\newcounter{utterance}

\centering \large Interaction Transcript for game `bbh', experiment `penguins\_in\_a\_table', episode 22 with reflexion{-}llama{-}3.1{-}8b{-}instruct.
\vspace{24pt}

{ \footnotesize  \setcounter{utterance}{1}
\setlength{\tabcolsep}{0pt}
\begin{supertabular}{c@{$\;$}|p{.15\linewidth}@{}p{.15\linewidth}p{.15\linewidth}p{.15\linewidth}p{.15\linewidth}p{.15\linewidth}}
    \# & \multicolumn{2}{c}{Player} && \multicolumn{2}{c}{Game Master} \\
    \hline

    \theutterance \stepcounter{utterance}  
    & & & \multicolumn{4}{p{0.6\linewidth}}{
        \cellcolor[rgb]{0.9,0.9,0.9}{
            \makecell[{{p{\linewidth}}}]{
                \texttt{\tiny{[P1$\langle$GM]}}
                \texttt{You have attempted this same task before and did not succeed. These are the plans you wrote after those attempts.} \\
\texttt{Trial \#1: New plan: } \\
\\ 
\texttt{In my last attempt, I failed to identify the correct number of penguins more than 8 years old. I incorrectly included Gwen in the list of penguins more than 8 years old, when in fact her age is 8 years. Next time, I will carefully examine the ages of the penguins and only include those that are strictly greater than 8 years old. I will also make sure to exclude penguins with an age of 8 years. To do this, I will create a list of penguins with ages greater than 8 and then count the number of penguins in this list.} \\
\\ 
\texttt{Trial \#2: New plan: In my last attempt, I failed to identify the correct number of penguins more than 8 years old. I incorrectly included Gwen in the list of penguins more than 8 years old, when in fact her age is 8 years. Next time, I will carefully examine the ages of the penguins and only include those that are strictly greater than 8 years old. I will also make sure to exclude penguins with an age of 8 years. To do this, I will create a list of penguins with ages greater than 8 and then count the number of penguins in this list. Additionally, I will also make sure to check the options provided to see if any of them match the number of penguins I have in my list, to avoid overcounting or undercounting.} \\
\\ 
\texttt{Trial \#3: In my last attempt, I failed to identify the correct number of penguins more than 8 years old. I incorrectly included Gwen in the list of penguins more than 8 years old, when in fact her age is 8 years. Next time, I will carefully examine the ages of the penguins and only include those that are strictly greater than 8 years old. I will also make sure to exclude penguins with an age of 8 years. To do this, I will create a list of penguins with ages greater than 8 and then count the number of penguins in this list. Additionally, I will also make sure to check the options provided to see if any of them match the number of penguins I have in my list, to avoid overcounting or undercounting.} \\
\\ 
\texttt{New plan: In my last attempt, I failed to consider the possibility that the correct answer might not be among the options provided. Next time, I will carefully examine the options and consider the possibility that the correct answer might be outside of the provided options. I will also make sure to double{-}check my calculations and consider alternative scenarios to ensure that I arrive at the correct answer.} \\
\\ 
\\ 
\texttt{Answer questions about a table of penguins and their attributes.} \\
\\ 
\texttt{Q: Here is a table where the first line is a header and each subsequent line is a penguin:  name, age, height (cm), weight (kg) Louis, 7, 50, 11 Bernard, 5, 80, 13 Vincent, 9, 60, 11 Gwen, 8, 70, 15  For example: the age of Louis is 7, the weight of Gwen is 15 kg, the height of Bernard is 80 cm.  We now add a penguin to the table:} \\
\texttt{James, 12, 90, 12} \\
\texttt{How many penguins are less than 8 years old?} \\
\texttt{Options:} \\
\texttt{(A) 1} \\
\texttt{(B) 2} \\
\texttt{(C) 3} \\
\texttt{(D) 4} \\
\texttt{(E) 5} \\
\texttt{A: (B)} \\
\\ 
\texttt{Q: Here is a table where the first line is a header and each subsequent line is a penguin:  name, age, height (cm), weight (kg) Louis, 7, 50, 11 Bernard, 5, 80, 13 Vincent, 9, 60, 11 Gwen, 8, 70, 15  For example: the age of Louis is 7, the weight of Gwen is 15 kg, the height of Bernard is 80 cm.  Which is the youngest penguin?} \\
\texttt{Options:} \\
\texttt{(A) Louis} \\
\texttt{(B) Bernard} \\
\texttt{(C) Vincent} \\
\texttt{(D) Gwen} \\
\texttt{(E) James} \\
\texttt{A: (B)} \\
\\ 
\texttt{Q: Here is a table where the first line is a header and each subsequent line is a penguin:  name, age, height (cm), weight (kg) Louis, 7, 50, 11 Bernard, 5, 80, 13 Vincent, 9, 60, 11 Gwen, 8, 70, 15  For example: the age of Louis is 7, the weight of Gwen is 15 kg, the height of Bernard is 80 cm.  What is the name of the second penguin sorted by alphabetic order?} \\
\texttt{Options:} \\
\texttt{(A) Louis} \\
\texttt{(B) Bernard} \\
\texttt{(C) Vincent} \\
\texttt{(D) Gwen} \\
\texttt{(E) James} \\
\texttt{A: (D)} \\
\\ 
\texttt{Q: Here is a table where the first line is a header and each subsequent line is a penguin:  name, age, height (cm), weight (kg) Louis, 7, 50, 11 Bernard, 5, 80, 13 Vincent, 9, 60, 11 Gwen, 8, 70, 15  For example: the age of Louis is 7, the weight of Gwen is 15 kg, the height of Bernard is 80 cm.  How many penguins are more than 8 years old?} \\
\texttt{Options:} \\
\texttt{(A) 1} \\
\texttt{(B) 2} \\
\texttt{(C) 3} \\
\texttt{(D) 4} \\
\texttt{(E) 5} \\
\texttt{A:} \\
            }
        }
    }
     \\ \\

    \theutterance \stepcounter{utterance}  
    & \multicolumn{4}{p{0.6\linewidth}}{
        \cellcolor[rgb]{0.561,0.737,0.561}{
            \makecell[{{p{\linewidth}}}]{
                \texttt{\tiny{[P1$\rangle$GM]}}
                \texttt{To answer the question, we need to carefully examine the ages of the penguins and only include those that are strictly greater than 8 years old. We will create a list of penguins with ages greater than 8 and then count the number of penguins in this list.} \\
\\ 
\texttt{From the table, we can see that the ages of the penguins are:} \\
\texttt{{-} Louis: 7} \\
\texttt{{-} Bernard: 5} \\
\texttt{{-} Vincent: 9} \\
\texttt{{-} Gwen: 8} \\
\texttt{{-} James: 12} \\
\\ 
\texttt{We will create a list of penguins with ages greater than 8:} \\
\texttt{{-} Vincent: 9} \\
\texttt{{-} James: 12} \\
\\ 
\texttt{There are 2 penguins in this list, so the correct answer is:} \\
\\ 
\texttt{(B) 2} \\
            }
        }
    }
    & & \\ \\

    \theutterance \stepcounter{utterance}  
    & & & \multicolumn{2}{p{0.3\linewidth}}{
        \cellcolor[rgb]{0.9,0.9,0.9}{
            \makecell[{{p{\linewidth}}}]{
                \texttt{\tiny{[GM$|$GM]}}
                \texttt{Target: (A)} \\
            }
        }
    }
    & & \\ \\

    \theutterance \stepcounter{utterance}  
    & & & \multicolumn{2}{p{0.3\linewidth}}{
        \cellcolor[rgb]{0.9,0.9,0.9}{
            \makecell[{{p{\linewidth}}}]{
                \texttt{\tiny{[GM$|$GM]}}
                \texttt{game\_result = LOSE} \\
            }
        }
    }
    & & \\ \\

    \theutterance \stepcounter{utterance}  
    & & \multicolumn{4}{p{0.6\linewidth}}{
        \cellcolor[rgb]{0.95,0.95,0.95}{
            \makecell[{{p{\linewidth}}}]{
                \texttt{\tiny{[P1$\leftrightarrow$P]}}
                \texttt{Post{-}episode Reflexion (not sent to GM):} \\
\texttt{New plan: In my last attempt, I failed to consider the possibility that the correct answer might not be among the options provided. Next time, I will carefully examine the options and consider the possibility that the correct answer might be outside of the provided options. I will also make sure to double{-}check my calculations and consider alternative scenarios to ensure that I arrive at the correct answer.} \\
            }
        }
    }
    & \\ \\

\end{supertabular}
}

\end{document}
